% Paragraphes « Protocole d'évaluation » et « Limites » (BoT-IoT), rédigés le 2026-10-10.
% Sources des chiffres : quickstart-pytorch/results/fp_reduction/ (check_protocol.json,
% leak_check_*.json, test_duplicates_botiot.json, */test_metrics.json, */shap_dns.json)
% et quickstart-pytorch/lab_test/results/lab_detection_summary.txt.
% Le tableau référencé est dans comparison_test.tex (\label{tab:botiot-fp-variantes}).

\paragraph{Protocole d'évaluation.}
Les 1\,011\,136 flux de BoT-IoT, dont seulement 9\,543 flux normaux (0,94~\%), sont répartis
entre sept clients : après un mélange aléatoire, les données sont découpées en sept tranches,
et chaque client sépare sa tranche en 80~\% pour l'entraînement et 20~\% pour son test local.
Notre première évaluation reposait sur un jeu de test « centralisé » tiré indépendamment de ce
découpage. Une vérification a posteriori a montré que 79,8~\% de ses lignes faisaient partie
des données d'entraînement d'un client, et que le seuil de décision (0,09) avait été choisi en
balayant ce même jeu. Le jeu de test de TON-IoT présente la même construction (80,0~\% de
recouvrement). L'effet mesuré sur les scores reste faible : sur BoT-IoT, le F1-macro vaut
0,911 sur les lignes déjà vues contre 0,905 sur les lignes jamais vues, et sur TON-IoT 0,649
contre 0,646. Ce protocole est néanmoins méthodologiquement invalide, et nous l'avons
remplacé. Le nouveau jeu d'évaluation est l'union des tests locaux des clients, soit
202\,230 flux qu'aucun client n'a vus pendant l'entraînement. Nous l'avons partagé en deux
moitiés stratifiées par sous-catégorie d'attaque : une \emph{validation} (101\,115 flux, dont
954 normaux) et un \emph{test} (101\,115 flux, dont 955 normaux). La normalisation min-max est
ajustée uniquement sur les données d'entraînement. Tous les choix (seuil de décision, variante
de modèle, pondération des classes) sont faits sur la validation. Nous imposons de plus que le
seuil tombe au milieu d'un écart d'au moins $10^{-3}$ entre deux scores consécutifs, pour ne
pas couper au milieu d'un amas de flux quasi identiques. Le seuil retenu, 0,167, correspond à
un taux de faux positifs inférieur à 5~\% sur la validation. Il a été figé, avec la variante
choisie, avant l'unique évaluation sur le test (tableau~\ref{tab:botiot-fp-variantes}). Sur ce
test, le modèle détecte 97,7~\% des attaques avec 4,5~\% de faux positifs (IC à 95~\% :
3,4--6,0~\%), contre 18,4~\% au seuil d'origine. L'AUC vaut 0,996 et la \emph{balanced
accuracy} 0,966. Le F1-macro (0,711) et le MCC (0,511) sont plus bas qu'au seuil d'origine, car
avec 99~\% d'attaques, les 2,3~\% d'attaques manquées représentent plus de lignes que tous
les flux normaux. Enfin, 20~\% des flux de test ont un contenu identique à un flux
d'entraînement, ce qui reflète le caractère répétitif du trafic d'attaque. Sur les seuls flux
uniques, les résultats sont équivalents (4,9~\% de faux positifs, 98,7~\% de rappel) : ils ne
sont donc pas gonflés par ces doublons.

\paragraph{Limites.}
Trois limites nuancent ces résultats. Premièrement, la séparation des flux normaux de petite
taille repose en partie sur un artefact du jeu de données. Au seuil d'origine, 71~\% des faux
positifs de validation (140 sur 196) sont des flux de type DNS : UDP, état INT, au plus quatre
paquets. Ils obtiennent un score d'environ 0,116, dans le même amas que les quelque 10\,300
attaques DoS et DDoS UDP de profil identique. L'analyse SHAP montre que l'état \texttt{INT}
pousse ces flux vers la classe « attaque », et que c'est la variable \texttt{hour\_of\_day}
(contribution moyenne $-0{,}29$) qui les retient du côté normal. Or cette variable traduit l'heure
de capture du jeu de données, et non une propriété du trafic. Un modèle réentraîné sans elle
(même configuration fédérée, 100 rounds) classe, au seuil 0,09, 329 de ces 336 flux normaux comme attaques. À
taux de faux positifs comparable (environ 4,4~\%), son rappel tombe à 95,0~\%, contre 97,7~\%.
Une pondération renforcée de la classe normale (poids multiplié par cinq) n'améliore pas cette
séparation : elle ne fait que décaler les scores. L'effet de l'heure se mesure directement sur le
trafic du laboratoire : pour un même flux normal, en ne faisant varier que \texttt{hour\_of\_day},
la probabilité d'attaque médiane passe d'environ 0,48 à minuit à 0,98 à 14~heures. Fixer cette
variable à sa valeur médiane d'entraînement, plutôt que de laisser l'heure réelle de capture, est
donc nécessaire pour que le seuil ne dépende pas de l'heure. Une analyse des variables disponibles
montre toutefois qu'une grandeur purement comportementale et reproductible, le nombre moyen
d'octets par paquet (\texttt{bytes}/\texttt{pkts}), sépare presque parfaitement ces deux groupes
sur la validation (aire sous la courbe de 0,99 : les attaques UDP se concentrent autour de
57~octets par paquet, les flux DNS normaux autour de 88). Un constat annexe mérite d'être relevé :
dans BoT-IoT, \emph{aucun flux UDP n'a de réponse capturée} --- les octets et paquets en retour
(\texttt{dbytes}, \texttt{dpkts}) sont nuls aussi bien pour les flux DNS normaux que pour les
attaques. La présence d'une réponse, qui distinguerait une requête DNS légitime d'un flot d'envoi,
n'est donc pas exploitable sur ce jeu. Par ailleurs, les ports ne figurent pas parmi les
variables : le terme « DNS » désigne un profil de flux, sans vérification du port~53.
Deuxièmement, certaines catégories sont trop rares pour être évaluées de façon fiable. Le vol de
données ne compte que 118 flux dans tout le jeu de données, dont 11 dans le test. Son rappel au
seuil retenu (4 sur 11, soit 36~\%, avec un IC à 95~\% de 15 à 65~\%) n'est pas interprétable,
et le keylogging (152 flux de test) reste lui aussi peu représenté. Troisièmement, les
performances mesurées sur BoT-IoT ne se transfèrent pas telles quelles à du trafic réel. Nous
avons testé le modèle dans un laboratoire isolé (victime Debian, attaquant Kali ; balayage nmap,
SYN flood avec hping3, force brute SSH), avec des variables extraites par tshark, sur 66\,588
flux. Les flux sont étiquetés par l'adresse IP de l'attaquant, réservée aux attaques pendant la
capture d'attaque. Un premier étiquetage par fenêtres de temps avait classé à tort le balayage
nmap (1\,000 flux, commencé 0,6~s après la fin de sa fenêtre) comme trafic normal. Il reste
seulement 49 flux normaux, capturés à part (SSH, HTTP et ICMP entre les deux machines), pour
66\,539 flux d'attaque. L'AUC vaut 0,99, mais sur 49 flux normaux seulement, et les scores sont
décalés : tous les flux normaux obtiennent une probabilité d'attaque comprise entre 0,977
et 1. On observe donc 100~\% de faux positifs aux seuils appris sur BoT-IoT (0,167 et 0,09).
Autrement dit, le modèle classe les flux normaux sous les attaques, mais aucun seuil appris
sur BoT-IoT ne les laisse passer. Ce décalage de domaine a deux causes. D'une part, les
variables les plus influentes du modèle ne sont pas reproduites à l'identique hors de l'outil
Argus utilisé par les auteurs du jeu de données. C'est le cas en particulier de la moyenne, de
l'écart-type, du minimum et du maximum, calculés sur des enregistrements agrégés selon des
paramètres non publiés. D'autre part, avec 99~\% d'attaques dans les données d'entraînement,
le modèle tend vers la classe majoritaire dès qu'il sort de sa distribution. Vu le faible
nombre de flux normaux, issus d'un seul couple de machines, cette AUC reste fragile. Un
déploiement réel exigerait au minimum de calibrer le seuil sur le trafic local, et de
préférence de réentraîner le modèle avec le même extracteur de variables que celui utilisé en
production.
